%% Lexical Structure
\chapter{Lexical Structure}
\label{ch:lexical}

The lexical structure of Lambdust extends R7RS-small with additional token types to support gradual dependent typing, effects, concurrency, and other advanced features while maintaining complete backward compatibility.

\section{Lexical Elements}
\label{sec:lexical-elements}

A Lambdust program consists of a sequence of lexical elements:
\begin{itemize}
\item Whitespace and comments (ignored)
\item Tokens (significant elements)
\item End-of-file marker
\end{itemize}

\section{Character Set}
\label{sec:charset}

Lambdust source code uses Unicode UTF-8 encoding. All R7RS-small character classifications are preserved:

\begin{itemize}
\item \textbf{Whitespace}: Space, tab, newline, carriage return, form feed
\item \textbf{Delimiters}: Whitespace, parentheses, brackets, braces, quotes, semicolon
\item \textbf{Letters}: Unicode letter categories (Alphabetic property)  
\item \textbf{Digits}: ASCII digits 0-9
\item \textbf{Special Initial}: \code{! \$ \% \& * / : < = > ? \^ \_ \textasciitilde}
\item \textbf{Special Subsequent}: \code{+ - . @}
\end{itemize}

\begin{compatibility}[Character Set Compatibility]
All R7RS-small character classifications and identifier rules are preserved unchanged in Lambdust.
\end{compatibility}

\section{Comments}
\label{sec:comments}

Lambdust supports all R7RS comment forms plus nested block comments:

\subsection{Line Comments}
\label{subsec:line-comments}

Line comments begin with semicolon and extend to end of line:
\begin{lambdustcode}
; This is a line comment
(define x 42) ; Another comment
\end{lambdustcode}

\subsection{Block Comments}  
\label{subsec:block-comments}

Block comments are delimited by \code{\#|} and \code{|\#} and may be nested:
\begin{lambdustcode}
#| This is a block comment
   #| with nested comments |#
   that can span multiple lines |#
(define y 17)
\end{lambdustcode}

\subsection{Expression Comments}
\label{subsec:expr-comments}

Expression comments use \code{\#;} to comment out the following expression:
\begin{lambdustcode}
(list 1 #;(* 2 3) 4)  ; Equivalent to (list 1 4)
\end{lambdustcode}

\section{Identifiers}
\label{sec:identifiers}

\indexentry{identifier}
Identifiers follow R7RS rules with no extensions:

\begin{grammar}
\production{\var{identifier}}{\produces}{\var{initial} \arbno{\var{subsequent}}}{}
\production{}{\alternative}{\var{vertical-line} \arbno{\var{symbol-element}} \var{vertical-line}}{}
\production{}{\alternative}{\var{peculiar-identifier}}{}
\\
\production{\var{initial}}{\produces}{\var{letter}}{}
\production{}{\alternative}{\var{special-initial}}{}
\\
\production{\var{letter}}{\produces}{a \alternative{} b \alternative{} c \alternative{} \ldots \alternative{} z}{}
\production{}{\alternative}{A \alternative{} B \alternative{} C \alternative{} \ldots \alternative{} Z}{}
\\
\production{\var{special-initial}}{\produces}{! \alternative{} \$ \alternative{} \% \alternative{} \& \alternative{} *}{}
\production{}{\alternative}{/ \alternative{} : \alternative{} < \alternative{} = \alternative{} >}{}
\production{}{\alternative}{? \alternative{} \^ \alternative{} \_ \alternative{} \~}{}
\\
\production{\var{subsequent}}{\produces}{\var{initial} \alternative{} \var{digit}}{}
\production{}{\alternative}{\var{special-subsequent}}{}
\\
\production{\var{special-subsequent}}{\produces}{+ \alternative{} - \alternative{} . \alternative{} @}{}
\\
\production{\var{peculiar-identifier}}{\produces}{+ \alternative{} - \alternative{} ... \alternative{} \goesto}{}
\end{grammar}

\begin{example}[Valid Identifiers]
\begin{lambdustcode}
lambda          ; Standard identifier
list->vector    ; With special characters
string-length   ; Hyphenated style
camelCase       ; Camel case style  
PascalCase      ; Pascal case style
+               ; Single character
...             ; Ellipsis
->              ; Arrow (peculiar identifier)
\end{lambdustcode}
\end{example}

\section{Keywords}
\label{sec:keywords}

\indexentry{keyword}
Lambdust extends R7RS with keyword literals for named parameters and metadata:

\begin{grammar}
\production{\var{keyword}}{\produces}{\#: \var{identifier}}{}
\end{grammar}

Keywords are self-evaluating and distinct from symbols:

\begin{lambdustcode}
#:type        ; Type annotation keyword
#:inline      ; Optimization hint keyword  
#:doc         ; Documentation keyword
#:test        ; Test specification keyword
\end{lambdustcode}

\begin{example}[Keyword Usage]
\begin{lambdustcode}
(define (factorial n)
  #:type (-> Natural Natural)
  #:doc "Compute factorial of natural number"
  (if (<= n 1) 1 (* n (factorial (- n 1)))))
\end{lambdustcode}
\end{example}

\section{Numbers}
\label{sec:numbers}

\indexentry{number}
Lambdust supports the complete R7RS numeric tower with all exactness and radix specifications:

\begin{grammar}
\production{\var{number}}{\produces}{\var{num-2} \alternative{} \var{num-8} \alternative{} \var{num-10} \alternative{} \var{num-16}}{}
\end{grammar}

\subsection{Integer Numbers}
\label{subsec:integers}

\begin{lambdustcode}
42          ; Decimal integer
#b1010      ; Binary integer (10)
#o52        ; Octal integer (42)  
#x2A        ; Hexadecimal integer (42)
#e42        ; Exact integer
#i42        ; Inexact integer
+42         ; Positive integer
-42         ; Negative integer
\end{lambdustcode}

\subsection{Real Numbers}
\label{subsec:reals}

\begin{lambdustcode}
3.14159     ; Decimal real
1.23e10     ; Scientific notation
1.23E-5     ; Scientific notation  
#i3.14      ; Explicitly inexact
#e1.5       ; Exact rational (3/2)
.5          ; Leading decimal point
5.          ; Trailing decimal point
\end{lambdustcode}

\subsection{Rational Numbers}
\label{subsec:rationals}

\begin{lambdustcode}
1/2         ; One half
-3/4        ; Negative three quarters
#e22/7      ; Exact rational
#i1/3       ; Inexact rational
\end{lambdustcode}

\subsection{Complex Numbers}
\label{subsec:complex}

\begin{lambdustcode}
3+4i        ; Rectangular form
3-4i        ; Rectangular form
+i          ; Pure imaginary unit
-i          ; Negative imaginary unit
3.0+0i      ; Real with imaginary zero
0+1.5i      ; Pure imaginary
\end{lambdustcode}

\section{Strings}
\label{sec:strings}

\indexentry{string}
String literals follow R7RS syntax with full escape sequence support:

\begin{grammar}
\production{\var{string}}{\produces}{" \arbno{\var{string-element}} "}{}
\\
\production{\var{string-element}}{\produces}{\var{any character other than " or \textbackslash}}{}
\production{}{\alternative}{\textbackslash{} " \alternative{} \textbackslash{} \textbackslash}{}
\production{}{\alternative}{\textbackslash{} n \alternative{} \textbackslash{} r \alternative{} \textbackslash{} t}{}
\production{}{\alternative}{\textbackslash{} \var{intraline-whitespace} \var{newline} \var{intraline-whitespace}}{}
\production{}{\alternative}{\textbackslash{} x \var{hex-digit} \arbno{\var{hex-digit}} ;}{}
\end{grammar}

\begin{example}[String Literals]
\begin{lambdustcode}
"Hello, World!"           ; Simple string
"String with \"quotes\""  ; Escaped quotes
"Line 1\nLine 2"         ; Newline escape
"Unicode: \x03BB;"       ; Unicode lambda
"Long string \
can span lines"          ; Line continuation
""                       ; Empty string
\end{lambdustcode}
\end{example}

\section{Characters}
\label{sec:characters}

\indexentry{character}
Character literals use R7RS syntax with named characters:

\begin{grammar}
\production{\var{character}}{\produces}{\#\textbackslash{} \var{any-character}}{}
\production{}{\alternative}{\#\textbackslash{} \var{character-name}}{}
\production{}{\alternative}{\#\textbackslash{} x \var{hex-scalar-value}}{}
\end{grammar}

\begin{example}[Character Literals]
\begin{lambdustcode}
#\a           ; Letter 'a'
#\A           ; Letter 'A'  
#\space       ; Space character
#\newline     ; Newline character
#\tab         ; Tab character
#\return      ; Carriage return
#\x03BB       ; Unicode lambda ($\lambda$)
#\x0          ; Null character
\end{lambdustcode}
\end{example}

\section{Booleans}
\label{sec:booleans}

\indexentry{boolean}
Boolean literals follow R7RS with both short and long forms:

\begin{grammar}
\production{\var{boolean}}{\produces}{\#t \alternative{} \#f \alternative{} \#true \alternative{} \#false}{}
\end{grammar}

\begin{lambdustcode}
#t      ; True (short form)
#f      ; False (short form)
#true   ; True (long form)  
#false  ; False (long form)
\end{lambdustcode}

\section{Delimiters}
\label{sec:delimiters}

\indexentry{delimiter}
Lambdust uses R7RS delimiters plus additional ones for type expressions:

\subsection{Parentheses}
\label{subsec:parens}

\begin{lambdustcode}
( )     ; Standard list delimiters
\end{lambdustcode}

\subsection{Brackets}
\label{subsec:brackets}

\begin{lambdustcode}
[ ]     ; Alternative list delimiters (equivalent to parentheses)
\end{lambdustcode}

\subsection{Braces}
\label{subsec:braces}

\begin{lambdustcode}
{ }     ; Set literals and type expressions
\end{lambdustcode}

\section{Quote and Unquote}
\label{sec:quote}

\indexentry{quote}
\indexentry{quasiquote}
\indexentry{unquote}
All R7RS quote forms are supported:

\begin{grammar}
\production{\var{abbreviation}}{\produces}{\var{abbrev-prefix} \var{datum}}{}
\\
\production{\var{abbrev-prefix}}{\produces}{\textquotesingle{} \alternative{} \textasciigrave{} \alternative{} , \alternative{} ,@}{}
\end{grammar}

\begin{example}[Quote Forms]
\begin{lambdustcode}
'(a b c)        ; Quote
`(list ,x ,@y)  ; Quasiquote with unquote and unquote-splicing
\end{lambdustcode}
\end{example}

\section{Lambdust Extensions}
\label{sec:extensions}

\indexentry{type annotation}
\indexentry{arrow}
\indexentry{pipe}
\indexentry{tilde}

Lambdust adds several new token types for advanced features:

\subsection{Type Annotation}
\label{subsec:type-annotation}

\begin{grammar}
\production{\var{type-annotation}}{\produces}{::}{}
\end{grammar}

Used for explicit type ascription:
\begin{lambdustcode}
(:: expr Type)              ; Type ascription
(define foo :: Type expr)   ; Definition with type
\end{lambdustcode}

\subsection{Arrow Types}
\label{subsec:arrows}

\begin{grammar}
\production{\var{arrow}}{\produces}{\code{->}}{}
\production{\var{fat-arrow}}{\produces}{=>}{}  
\production{\var{tilde-arrow}}{\produces}{\~>}{}
\end{grammar}

Used in type expressions:
\begin{lambdustcode}
(-> A B)        ; Function type A to B
(=> P Q)        ; Implication or constraint  
(~> A B)        ; Effectful function type
\end{lambdustcode}

\subsection{Type Operators}
\label{subsec:type-ops}

\begin{grammar}
\production{\var{colon}}{\produces}{:}{}
\production{\var{pipe}}{\produces}{|}{}
\production{\var{tilde}}{\produces}{~}{}
\end{grammar}

Used in type expressions:
\begin{lambdustcode}
(x : Type)      ; Parameter with type
A | B           ; Union type or type alternative
~Effect         ; Effect annotation
\end{lambdustcode}

\section{Token Classification}
\label{sec:token-classification}

The lexer recognizes the following token types:

\onecolumn
\begin{table}[H]
\centering
\caption{Lambdust Token Types}
\label{tab:tokens}
\begin{tabular}{|l|l|l|}
\hline
\textbf{Token Type} & \textbf{Examples} & \textbf{Usage} \\
\hline
LeftParen & \code{(} & List start, procedure calls \\
RightParen & \code{)} & List end, procedure calls \\
LeftBracket & \code{[} & Alternative list notation \\
RightBracket & \code{]} & Alternative list notation \\
LeftBrace & \code{\{} & Set literals, type expressions \\
RightBrace & \code{\}} & Set literals, type expressions \\
Quote & \code{'} & Quote abbreviation \\
Quasiquote & \code{`} & Quasiquote abbreviation \\
Unquote & \code{,} & Unquote abbreviation \\
UnquoteSplicing & \code{,@} & Unquote-splicing abbreviation \\
Dot & \code{.} & Pair notation, module access \\
TypeAnnotation & \code{::} & Type ascription \\
Colon & \code{:} & Parameter type annotation \\
Arrow & \code{->} & Function types \\
FatArrow & \code{=>} & Constraints, implications \\
Pipe & \code{|} & Union types, alternatives \\
Tilde & \texttt{\textasciitilde} & Effect annotations \\
TildeArrow & \texttt{\textasciitilde>} & Effectful function types \\
IntegerNumber & \code{42}, \code{\#x2A} & Integer literals \\
RealNumber & \code{3.14}, \code{1e10} & Real number literals \\
RationalNumber & \code{1/2}, \code{22/7} & Rational literals \\
ComplexNumber & \code{3+4i}, \code{+i} & Complex number literals \\
String & \code{"hello"} & String literals \\
Character & \code{\#\textbackslash{}a}, \code{\#\textbackslash{}space} & Character literals \\
Boolean & \code{\#t}, \code{\#false} & Boolean literals \\
Keyword & \code{\#:type}, \code{\#:doc} & Keyword literals \\
Identifier & \code{lambda}, \code{+}, \code{foo} & Variable names, operators \\
LineComment & \code{; comment} & Documentation (ignored) \\
BlockComment & \code{\#| comment |\#} & Documentation (ignored) \\
Error & Invalid input & Lexical errors \\
Eof & End of input & End of file marker \\
\hline
\end{tabular}
\end{table}
\twocolumn

\section{Lexical Precedence}
\label{sec:precedence}

When multiple token patterns could match, the lexer uses longest-match with the following precedence:

\begin{enumerate}
\item Numbers (including complex forms)
\item Keywords (\code{\#:identifier})
\item Booleans (\code{\#t}, \code{\#f}, \code{\#true}, \code{\#false})
\item Characters (\code{\#\textbackslash{}...})
\item Multi-character operators (\code{::}, \code{->}, \code{=>}, \code{\~>}, \code{,@})
\item Single-character tokens
\item Identifiers
\item Comments
\end{enumerate}

\section{Error Handling}
\label{sec:lexical-errors}

\indexentry{lexical error}
The lexer reports errors for:

\begin{itemize}
\item Invalid character sequences
\item Unterminated strings or comments
\item Invalid numeric formats
\item Invalid character literals
\item Invalid Unicode escape sequences
\end{itemize}

Each error includes source location information (line, column, file) for precise error reporting.

\section{Implementation Notes}
\label{sec:lexical-implementation}

\subsection{Performance Considerations}
The lexer is optimized for:
\begin{itemize}
\item Single-pass tokenization
\item Minimal backtracking
\item Efficient Unicode handling
\item Memory-efficient token storage
\end{itemize}

\subsection{Compatibility Mode}
Implementations may provide a compatibility mode that accepts only R7RS tokens, rejecting Lambdust extensions for strict R7RS compliance checking.

\subsection{Extensibility}
The lexical structure is designed to allow future extensions without breaking existing code. Reserved token patterns may be allocated for future language features.